\chapter{Summary and Outlook}
\label{summary}

With each section of this thesis now complete it is time to take
a step back and consider the work from a broader perspective,
reflecting on what was achieved, how it was achieved,
and what this means for the future.

First the work done in this thesis will be summarized on a high level.
Then the questions raised in the introduction will be answered and
the contributions of this thesis will be outlined.
Afterwards the work done, and the choices made during this project will
be reflect on.
Finally, the outlook will be given, discussing the steps following
the completion of this thesis and the future work to be done.

\section{Summary}


%  rehashin the goals, the problem and the contributions we aim to make



\textbf{Introduction:}This research project was inspired by the need for a user-friendly way to interact with
\ac{HPC} systems, which themselves should be easy and quick to create.
The need for an interactive environment being driven directly by the shortcomings of an earlier project.
This being a pipeline driven \gls{GitOps} system which automatically schedules \ac{HPC} workloads on a Kubernetes cluster.
As the jobs are scheduled by the pipeline, which is being build by committing to a git repository,
this system lacked the interactivity needed to develop the pipeline code.
For this reason the goal for this project was to create a user-friendly way,
to enable the user to interact with the \ac{HPC} system in a more direct
% to do the exploitative computing needed for the development of the pipeline code
% or one off jobs which do not need to be run in the pipeline.
Based on this problem the problem statement,
research questions and the contributions of this thesis were defined.


\textbf{Methodology: }To solve this problem which exists in the intersection the \ac{HPC} and \ac{CC} worlds,
which are both highly complex and specialized fields,
a combination of methodologies was chosen which on one hand allows for a deep
understanding of the theoretical background and on the other hand contrasts
this with a practical implementation.

For this reason a conceptual deductive Analysis was chosen, in order
to derive a design from the theoretical background and the requirements from the problem statement.
This is done by starting as close to the first principles of the
technologies involved as possible and then working towards the higher levels required for the implementation.

To prove or invalidate this hypothesis of a system design it was
decided to implement the as a prototype in a series of iterations and
then evaluate the results of these iterations against the design and the theoretical background. \

% - theory
%  - how cloud works
%     how isolation works, 
%     - this is important to understand the context of container,
%     - the subject of kvm shows up again with \ac{SRIOV}

%     how this evolved into containers and how those work
%     - this is importat to under stand what to do in the protyping
%     - understanding the cgroups and capabiliets was also important for understanding 
%     how to set up the containers for SRIOV
%     - the knowledge gainded of container runtimes and images helped to during the second iteration

\pagebreak


\textbf{Theoretical Work - Cloud:} The theoretical work started with the foundations of the \ac{CC} world.
Beginning with the basics of process and systems isolation and how they are achieved in practice.
Starting at \ac{KVM} and their ability to access some of the hardware directly,
did not only provide insight into the predecessor of containers but also into the features
which became important later during the use of \ac{SRIOV} devices.
The inner workings of containers and their implementation in the Linux kernel were addressed next,
as well as how container runtime and their images work.
Bringing valuable insight especially useful during the second stage of the prototype.

% - how kubernetes works is of selv evident importance for the project
% how the networking does even more so
% but the scheduling processes where inteded for the 4th iteraton which did not come to pass

Expanding from containerization to the concepts of Kubernetes,which first where explained on an architectural level,
and then on a networking level Information which served as the basis for each of the iterations.
The final section, explaining the scheduling and topology considerations of Kubernetes where intended as basis
for the fourth iteration and included in the design of the system, but where not implemented due to time constraints.

\textbf{Theoretical Work - HPC:} Having established the theoretical background of the \ac{CC} world,
the background of the \ac{HPC} world was addressed next.
Once again starting of with the very basics, the distributed processing and relevant
hardware infrastructure, in order to establish the context for the relevance of Memory architectures.
As the problems of \ac{HPC} are much closer to the hardware than the typical \ac{CC} applications,
this knowledge is essential to understand the way these system allocate and access memory as this
builds the foundation of more complex memory architectures like the \ac{PGAS} model relying
on a \ac{NUMA} like architecture of over the network memory access.
Leading into the explanation of network interconnects and fabrics and the place
they hold through the \ac{RDMA} and \ac{MPI} protocols.
From there the complex network topologies, common in the \ac{HPC} world to reduce cost and
increase performance where explained.
Finally these concepts are tied together with an overview of the most relevant software stack components
from the \ac{OS} to \textit{Workload Management Systems} are presented.

% - how hpc works 
%     as hpc has been evolving in paralel to other compute technologuies,
%         it is important to understand the considerations and problems these systeems
%         face and solve. 
%     - Processing and the basics of distributed processing is explained
%     - then memory placement and and the problems of distributed memory are explained resulting in
%     an explanation of the \ac{PGAS} model which esenntialy relies on NUMA on steroids over the network
%     - the netwok interconnects are explained in detail, as they are one of the most important parts of an hpc system
%     - then Topologies are explained, which would have become important later
%     - then a short overview of the software stack is given, which is important to understand the context of the chapel language
%     - then workload management systesms are explained to contrast the kubernetes scheduler

\textbf{Theoretical Work - Exploratory Computing:} Last but not least some of the theoretical background of
the current state of the art in exploratory computing was presented.
A field which is highly influenced by the computational notebook paradigm,
which is a way to combine code, documentation and data visualization in a single document,
which has become very popular in the data science community and
many other fields where exploratory computing is done.
%     - a short overview of exploratory computing and notebooks is given


% - practice 
%     - desing of the system
%         - the problems from the begining are converted into a goal of a system
%         - based on the information gathered in the previous chapters the requirements to achieve this goal are derived
%         - upon a these requirements a design is constructed, going through each of the segments discussed in the theoretical chapters
%           and a design decision is made based on that information

\textbf{Experimental Work - Design:} With the theoretical background established the design of the system could be derived.
Starting with the problem statement from the introduction the goals for the design where derived.
These goals where then used to inform the requirements for the system and a design was constructed based on these requirements.
The design itself was broken into each of the segments discussed in the theoretical chapters and at each steps
a design decision was made based on the information gathered in the theoretical chapters.


% - aritfact creation 
%     - first the theroetical achritecture gets mapped to  the realities of this prototype and the previous work this builds on
%     - an evaluation of which elements of the previous work fit the constucted design and which do not is made
%     - the iterative development starts each iteration serving as a the basis for the next and informes the design of the next


\textbf{Experimental Work - Artifact Creation:} Before this design was implemented as a prototype,
the theoretically designed system which was targeted to be a general design with as little influence form previous work as possible,
was mapped to the practicalities of the environment the prototype was to be implemented in, discussing the previous work.
the existing infrastructure and how this translates to the design of the system to be implemented.

% - creating a full blown MVP of the system including fixes and the building, setup and integration of all new components
% - the quick and dirty MVP is then rebuild from scratch in order to create a more performant and resonsive system, 
%   additionally a performance meassurment is integrated and it is attempted to integrate the singularity \ac{CRI} which
%   only now reveals itself as as an impossibility
% - the third iteration then aims to integrate infiniband infrasturcture to eliminate the bottleneck of the network
%   integrating the whole networkstack using a fully automated configuration setup of the interface device driver and 
%   network stack results in running out of time before it could be concluded or tested 

\textbf{Experimental Work - Iterative Development:} With the actual environment established the iterative development of the prototype
was started. Each iteration setting a goal based on the design and the previous work, describing the implementation of these goals,
mentioning mayor problems and solutions found during the implementation and evaluating the results of each step.
The fist iteration started of by implementing a \ac{MVP} of the system, integrating the all the software components needed for the system to work.
During the second iteration the all containers of the system where rebuild from scratch to create a more performant and responsive
system, and converting the benchmarking system of Arkouda so it could be used to evaluate the performance of Kubernetes
cluster. Additionally it was attempted to integrate the \ac{CRI} of Singularity, which only then revealed itself as not feasible.
The third iteration aimed to replace the networking layer with a fully automated multi \ac{CNI} setup,
enabling the worker nodes to communicate more efficiently with the \ac{IB} network,
in order to eliminate the bottleneck discovered in the second iteration.
This last iteration, while in the final stages could not be completed due to time constraints.
The topics of dynamics autoscaling and cluster topography, while in the original design, where not implemented due the
same time constraints.

\textbf{Experimental Work - Evaluation:} With the development process concluded,
the final step of the deductive design process was to contrast the deductive design with the discoveries made during the
experimental work. This was done by going through each of the design decisions made in the theoretical design and
evaluating them separately against the results of the experimental work.
Validating, correcting or invalidating each of the design decisions based on the new information gathered during
the prototyping process and determining which changes made during the prototyping process would be necessary to
an unrelated implementation and which where only necessary due to the specific constraints of the prototype.
% - this brings and end to the iterative process, preventing the implementation of the other subjects set up in
%   theory and design.

% - the iterations are concluded with a summary of the process 
% - and once again each part of the system design is rehashed and the final system is used to map the decisions 
%     made in the theoretical architecture against a real world implementation
% - the whole 

\pagebreak

\subsection{Contributions}

With the work on this thesis now concluded,
the contribution of this thesis can be summarized and the research questions answered.

% - answering the research questions
\subsubsection{Research Questions}
First of lets summarize the answer this thesis has given to the research questions:

\begin{itemize}
    \item \textbf{RQ1:} \textit{
              How can the system be extended to allow for a real-time explorative workflow, allowing for an interactive exploration of \ac{HPC} scale data?
          }
    \item \textbf{Answer:} Within this thesis an interactive interface for an Arkouda based Compute Cluster was designed and created as an artifact.
          By way of confirming and correcting the original deductive design with the knowledge gained during an experimental implementation,
          a new design was created which describes how such a system can be implemented in detail. The main components being a
          computational notebook interfacing with a worker fleet of Arkouda workers running on a Kubernetes cluster, communicating over
          \ac{PGAS} and if possible utilizing \ac{RDMA} protocols.
    \item \textbf{RQ2:} \textit{
              How can such a system be designed to combine the performance of \ac{HPC} with the flexibility and convenience of \ac{CC} systems?
          }
    \item \textbf{Answer:} During the theoretical and experimental work the set of design decisions and tradeoffs needed to move towards such a
          system where described.
          The most significant factors being that while containerized systems have near native performance in computation, they suffer form
          a highly complex stack. In order to reach the performance of \ac{HPC} systems the network stack needs to be replaced
          with a more performant one, creating a highly complex system in order to be able to use the \ac{HPC} interconnect in a flexible
          and convenient way. The main tradeoffs of such a system being the complexity of the system and the performance of it.
          For more details on the design decisions and tradeoffs necessary see the evaluation section.

\end{itemize}

\subsubsection{Artifacts}

The direct contributions to the larger field of \ac{HPC} and \ac{CC} are as follows:

\begin{itemize}
    \item \textbf{C1:} \textit{ A deductive analysis of the three intersecting fields of \ac{HPC}, \ac{CC} and \ac{EC} has been created and experimentally validated.}
    \item \textbf{C2:} \textit{ A design for a user friendly and scalable HPC environment on kubernetes has been created and experimentally validated.}
    \item \textbf{C3:} \textit{ A tangible artifact in form of an integration of Arkouda, Jupyter and Kubernetes has been created and made available as a software project.}
\end{itemize}

The work done in this thesis is a step towards a user friendly and scalable \ac{HPC} environment over all.
As this system integrates with the Pachykouda project, a significant progression towards an easy
to use and scale HPC environment has been made.
Allowing the a user friendly development of pipelines,directly developed on the same system as they will be run on,
and can be directly committed to the pipeline repository local to the system, also using the same data source and storage
as used by the pachyderm platform.

Some of the optimization of the system, like the optimization of the images will also be
useful for the Pachykouda project but might also be useful upstream in the Arkouda project.

\section{Reflection}


\subsection{Methodology}

The dual methodology of combining the \ac{CDA} with \ac{RAD} has proven to be a good choice for this project.
While the \ac{CDA} provided a deep understanding of the theoretical background and an academic approach to the problem,
the \ac{RAD} on the other hand allowed for a fast and engineering based approach to the problem, of creating an experimental system.
Both in conjunction provided a very useful synergy between them, making the project both deep and practical at the same time.
But the practical challenges encountered, particularly in integrating InfiniBand within Kubernetes,
also show that creating and estimating the complexity of such a system is not early done from literature alone.
The need to crate the theoretical underpinning before the practicalities of the system can be understood,
makes it easy to research and describe topics which will only be tangentially relevant to system eventually created.

\subsection{Execution}

The execution of the project was done by creating a theoretical foundation, a design and then implementing this design as a prototype.
During the theoretical work a large focus where put on the theoretical underpinnings of the highly complex
distribute systems and the \ac{EC} where reduced to a short description of the computational notebook paradigm.
The same can be said for the practical implementation, where the focus was on the integration of \ac{HPC} paradigms into a \ac{CC} system,
while low hanging fruit of the user facing aspects of the system where not explored in depth.
Integrations which could have provided important usability features, like a direct integration with the version control system or
the available block storage of the existing system in order to move away from a totally ephemeral environment where not explored,
even though the infrastructure was already in place to do so.

\subsection{Results}

In the end the crated artifact is fairly complete and fully usable.
The system is able to run Arkouda code on a Kubernetes cluster, with a Jupyter notebook interface,
and has been used personally to do data analysis on small data analysis tasks, by the author out
of convenience.
So while the sudden end of the project due to time constraints is unfortunate especially
as the last iteration was only one or two days away from completion, the results of the project are still a success.


\section{Outlook}

As the project itself has shown much potential and a solid foundation has been laid,
there are many ways to continue the work done in this thesis.
The most immediate step would be to offer the improvements of the components
created in this thesis to the Arkouda community.
% - based on the work already done a pull request to the arkodua contrip repository will be made in the near future
% - basis for a new arkouda-contrip repo, supporting ib-kubernetes

\subsection{Future Work}
\label{sec:future_work}

The most immediate steps which still needs to be worked on is to complete the last iteration of the project.
As the work is mostly done and the integration of the InfiniBand network stack has been of intersecting to
the Arkouda community before.
Based on this system the performance comparison between the baseline established during the second iteration
and the new system can be made, to see if the performance of the system has improved as expected.

In regards to iterations not started yet, the dynamic autoscaling of the system would be a logical and
interesting next step. Both as a good way to reduce cost and increase the peak performance of the system,
but also as a way to further explore the Kubernetes system and its capabilities out of personal interest.

Of a high academic ss well as practical and personal interest would be the use of an
externalized global memory bank for the Arkouda workers to use as a collaborative memory space.
A solution like a \ac{FAM} system has the potential to increase the performance beyond a fully
node based \ac{PGAS} system.

Additionally this project will be continued in the context of the Hewlett Packard Labs internships,
where aspects such as an optimized memory allocation strategy of the Arkouda workers will be explored
by fellow students furthering the research and impact of this project.

\subsection{Conclusion}

In conclusion there while the results of this project where a success,
there is still much that can be explored and investigated in the field of \ac{HPC} and \ac{CC} systems.
